% Charte graphique des figures K&N (style « D5 », choisi le 2026-11-13).
%
% Une figure TikZ écrite par un modèle ne contient que le dessin : tout ce qui
% fait l'uniformité (épaisseurs, couleurs, police, codages, cotes) vient d'ici.
% services/tikzfig.py refuse dans une figure les épaisseurs, polices, échelles
% et couleurs qui ne passent pas par cette charte.
%
% Toute modification de ce fichier change la clé des figures rendues
% (tikzfig.CHARTE_VERSION) : elles sont recompilées sur le Mac au prochain rendu.
\usepackage{fontspec}
\usepackage{unicode-math}
\setmainfont{FiraSans}[Path=@FONTS@/,Extension=.ttf,UprightFont=*-Regular,BoldFont=*-SemiBold]
\setmathfont{FiraMath-Regular.otf}
\usepackage{xcolor,etoolbox}
\usepackage{tikz,tikz-3dplot,tkz-euclide,pgfplots}
\pgfplotsset{compat=1.18}
\usetikzlibrary{calc,angles,quotes,arrows.meta,intersections,positioning,shapes.geometric,decorations.markings,decorations.pathreplacing,backgrounds,fit}

% ------------------------------------------------------------------ couleurs
% Trait, noms, cotes, accent : toujours ces couleurs-là.
\definecolor{ink}{HTML}{1F2D36}      % traits
\definecolor{nomc}{HTML}{0B3C8C}     % noms de points et de figures
\definecolor{cotec}{HTML}{C24100}    % chiffres des cotes, valeurs d'angle
\definecolor{accent}{HTML}{00838F}   % section, élément à mettre en valeur, codages
% Fonds : purement esthétiques, très clairs.
\definecolor{fillA}{HTML}{F0F4F7}
\definecolor{fillB}{HTML}{E3EAEF}
\definecolor{angfill}{HTML}{DDF0EE}
\definecolor{gridc}{HTML}{A9B7C4}   % grille des figures (jamais celle des tracés, saumon)
% Repérage : des couleurs franches, pour reconnaître une figure parmi d'autres
% (figure d'origine, images, solutions candidates). Les sujets sont imprimés en couleur.
\definecolor{repA}{HTML}{2F6FC0}     % bleu
\definecolor{repB}{HTML}{2E9E5B}     % vert
\definecolor{repC}{HTML}{8A4FC7}     % violet
\definecolor{repD}{HTML}{E07A1F}     % orange
\definecolor{repE}{HTML}{C83E7A}     % framboise

% ------------------------------------------------------------------ calques
% Ordre de dessin imposé, quel que soit l'ordre du code : les fonds (aplats, faces,
% sections) tout en bas, puis les codages d'angle, puis les traits et les textes.
\pgfdeclarelayer{fond}\pgfdeclarelayer{angles}\pgfsetlayers{fond,angles,main}
\tikzset{on layer/.code={\pgfonlayer{#1}\begingroup\aftergroup\endpgfonlayer\aftergroup\endgroup},
 fond/.style={fill=none,preaction={on layer=fond,#1}}}

% ------------------------------------------------------------------ styles
\tikzset{
 every picture/.style={line join=miter,line cap=butt,font=\fontsize{8.5}{10}\selectfont,color=ink},
 trait/.style={line width=0.6pt,draw=ink},
 fort/.style={line width=1.1pt,draw=accent},
 cache/.style={line width=0.35pt,draw=ink,dash pattern=on 3pt off 1.5pt},
 constr/.style={line width=0.3pt,draw=ink!55},
 aplat/.style={fond={fill=fillA}},
 aplatB/.style={fond={fill=fillB}},
 section/.style={fond={fill=repA!30},draw=repA!80!ink,line width=1.1pt},
 sectionbord/.style={draw=repA!80!ink,line width=1.1pt},
 codage/.style={line width=0.4pt,draw=accent},
 grille/.style={line width=0.4pt,draw=gridc},
 axe/.style={line width=0.6pt,draw=ink,-{Stealth[length=4pt,width=3pt]}},
 etiq/.style={inner sep=1.8pt,font=\fontsize{10}{11}\selectfont\bfseries,text=nomc},
 petit/.style={inner sep=1.2pt,font=\fontsize{8}{9}\selectfont,text=ink},
 cote/.style={font=\fontsize{9.5}{10}\selectfont\bfseries,inner sep=1.4pt,text=cotec,fill=white},
 cotetrait/.style={line width=0.3pt,draw=cotec!40},
 rappel/.style={line width=0.25pt,draw=ink!25},
 cotefl/.style={{Stealth[length=2.6pt,width=1.8pt]}-{Stealth[length=2.6pt,width=1.8pt]}},
 boite/.style={draw=ink,line width=0.6pt,fill=white,minimum height=8mm,minimum width=12mm,
   inner xsep=3mm,font=\fontsize{12}{14}\selectfont},
 % figures repérées : `origine` pour la figure de départ (aplat soutenu),
 % `figA`…`figE` (couleurs repA…repE) pour distinguer des figures entre elles
 origine/.style={fond={fill=repA!45},draw=ink,line width=0.7pt},
 figA/.style={fond={fill=repA!22},draw=repA!75!ink,line width=0.7pt},
 figB/.style={fond={fill=repB!22},draw=repB!75!ink,line width=0.7pt},
 figC/.style={fond={fill=repC!22},draw=repC!75!ink,line width=0.7pt},
 figD/.style={fond={fill=repD!22},draw=repD!75!ink,line width=0.7pt},
 figE/.style={fond={fill=repE!22},draw=repE!75!ink,line width=0.7pt},
 % solides : faces du dessus, de face, de côté ; surface de révolution
 faceH/.style={fond={fill=fillA!50}}, faceV/.style={fond={fill=fillA}}, faceD/.style={fond={fill=fillB}},
 cylL/.style={fond={left color=fillA,right color=fillB!85!ink!25,middle color=white}},
 cylT/.style={fond={fill=fillA!60}},
 pics/ptmark/.style={code={\draw[line width=0.45pt,draw=ink] (-1.5pt,-1.5pt)--(1.5pt,1.5pt) (-1.5pt,1.5pt)--(1.5pt,-1.5pt);}},
}
\def\AngR{5mm}

% ------------------------------------------------------------------ macros
% \Pt{(x,y)}{A}{above right} : point marqué d'une croix et nommé
\newcommand{\Pt}[3]{\path #1 pic{ptmark}; \node[etiq,#3] at #1 {#2};}
% \Nom{(x,y)}{A}{below} : nom seul, sans croix
\newcommand{\Nom}[3]{\node[etiq,#3] at #1 {#2};}
% \PtA{(A)}{A}{30} et \NomA{(A)}{A}{30} : nom posé à 3,2 mm dans la direction donnée
% (en degrés) ; choisir la direction la plus dégagée, à l'opposé des traits qui partent du point
\newcommand{\PtA}[3]{\path #1 pic{ptmark}; \node[etiq] at ($#1+(#3:3.2mm)$) {#2};}
\newcommand{\NomA}[3]{\node[etiq] at ($#1+(#3:3.2mm)$) {#2};}
% \AngDroit{A}{B}{C} : angle droit en B (points nommés par \coordinate)
\newcommand{\AngDroit}[3]{\path[on layer=angles] pic[codage,angle radius=3.2mm]{right angle=#1--#2--#3};}
% \AngMark{A}{B}{C}{valeur}{arcs} : angle en B, 1 ou 2 arcs, valeur ou {} ; 0 arc = valeur seule
\newcommand{\AngMark}[5]{%
  \ifnum#5>0 \path[on layer=angles] pic[draw=none,fill=angfill,angle radius=\AngR]{angle=#1--#2--#3};
    \path[on layer=angles] pic[codage,angle radius=\AngR]{angle=#1--#2--#3};\fi
  \ifnum#5>1 \path[on layer=angles] pic[codage,angle radius=\AngR+1.6pt]{angle=#1--#2--#3};\fi
  \path pic[draw=none,"{\textcolor{cotec}{\fontsize{9.5}{10}\selectfont\bfseries #4}}",
    angle eccentricity=1.8,angle radius=\AngR]{angle=#1--#2--#3};}
% \Cote{(A)}{(B)}{6 cm}{4mm} : cote fléchée parallèle à [AB], décalée de 4 mm
% à gauche de A→B (valeur négative : à droite), avec lignes de rappel
\newcommand{\Cote}[4]{%
  \draw[rappel] #1 -- ($#1!#4*1.25!90:#2$) #2 -- ($#2!#4*1.25!-90:#1$);
  \draw[cotetrait,cotefl] ($#1!#4!90:#2$) -- ($#2!#4!-90:#1$)
     node[cote,midway,sloped] {#3};}
% \CoteV{(x,y1)}{(x,y2)}{3 m} : cote verticale, texte droit à sa droite
\newcommand{\CoteV}[3]{%
  \draw[cotetrait,cotefl] #1 -- #2 node[cote,midway,right=1pt] {#3};}
% \Marque{(A)}{(B)}{n} : n traits de codage (longueurs égales) au milieu de [AB]
\newcommand{\Marque}[3]{\path #1 -- #2 node[midway,sloped,inner sep=0pt,text=accent,
  font=\fontsize{9}{9}\selectfont\bfseries] {\ifcase#3\or|\or||\or|||\fi};}
% \DroiteGraduee{min}{max} : droite graduée de min à max (entiers), une graduation par
% unité, chaque nombre écrit dessous. Les points se placent ensuite avec \Pt{(x,0)}{A}{above}.
\newcommand{\DroiteGraduee}[2]{%
  \draw[axe] (#1-0.5,0)--(#2+0.7,0);
  \foreach \v in {#1,...,#2}{\draw[trait] (\v,-2.5pt)--(\v,2.5pt);
    \node[petit,below=3pt] at (\v,0) {$\v$};}}
% \Balance[eq|gauche|droite]{plateau gauche}{plateau droit} : balance de Roberval.
% `gauche` = le plateau gauche est plus lourd (il descend).
\newcommand{\Balance}[3][eq]{%
  \def\kntilt{0}%
  \ifstrequal{#1}{gauche}{\def\kntilt{7}}{}%
  \ifstrequal{#1}{droite}{\def\kntilt{-7}}{}%
  \begin{scope}
  \filldraw[trait,aplatB] (-0.75,-1.05)--(0.75,-1.05)--(0.45,-0.85)--(-0.45,-0.85)--cycle;
  \filldraw[trait,aplat] (-0.07,-0.85) rectangle (0.07,0);
  \begin{scope}[rotate=\kntilt]
    \filldraw[trait,aplatB] (-2.1,-0.05) rectangle (2.1,0.05);
    \coordinate (knL) at (-1.85,0.05); \coordinate (knR) at (1.85,0.05);
  \end{scope}
  \filldraw[trait,fill=white] (0,0) circle (0.9mm);
  \foreach \s/\t/\c in {L/#2/knL,R/#3/knR}{
    \filldraw[trait,aplatB] ($(\c)+(-0.04,0)$) rectangle ($(\c)+(0.04,0.22)$);
    \node[boite,anchor=south] (knbox\s) at ($(\c)+(0,0.30)$) {\t};
    \filldraw[trait,aplatB] ($(knbox\s.south west)+(-0.18,-0.08)$) -- ($(knbox\s.south east)+(0.18,-0.08)$)
      -- ($(knbox\s.south east)+(0.08,0)$) -- ($(knbox\s.south west)+(-0.08,0)$) -- cycle;}
  \end{scope}}

% ------------------------------------------------------------------ pgfplots
\pgfplotsset{
 every axis/.append style={line width=0.6pt,axis line style={draw=ink},
   tick style={draw=ink,line width=0.4pt},
   tick label style={font=\fontsize{8}{9}\selectfont,text=ink},
   label style={font=\fontsize{9}{10}\selectfont,text=ink},
   major grid style={draw=gridc,line width=0.4pt},
   minor grid style={draw=gridc!55,line width=0.25pt},
   legend style={font=\fontsize{8}{9}\selectfont,draw=gridc}},
 every axis plot/.append style={line width=1pt},
 cycle list={{repA,mark=none},{repD,mark=none},{repB,mark=none},{repC,mark=none},{repE,mark=none}},
 kn repere/.style={axis lines=middle,grid=both,axis line style={-{Stealth[length=4pt,width=3pt]}},
   xlabel style={anchor=north west},ylabel style={anchor=south east}},
}
